% BEGIN S1
\par\medskip\noindent\textbf{1 Introduction}\par
% BEGIN S1.p1
The concept of distance balls is of fundamental importance to any given metric space $(X,d)$ as the concept of open sets to any topological space. A distance ball $B(v,r)$ in a metric space $(X,d)$ with radius $r$ centered at point $v$ is the subset of points $\{u\in X:d(u,v)\leq r\}$ in $X$. Many of the algorithmic success on graphs can be attributed to new understandings of the structure of distance balls in the corresponding graph metric.
\par
\par
% END S1.p1
% BEGIN S1.SS0.SSS0.P0.SPx1
\par\medskip\noindent\textbf{The role of distance balls.}\par
% BEGIN S1.SS0.SSS0.P0.SPx1.p1
Take the minimum dominating set – one of the 21 famous NP-complete problems of Karp – as an example. Given an unweighted graph $G$, the minimum dominating set problem (Dominating Set for short) asks for a subset of vertices $S$ of smallest cardinality, such that every vertex in $G$ is either in $S$, or is adjacent to some vertex in $S$. Dominating Set can also be interpreted as an instance of the set cover problem: Consider the set $\mathcal{R}\coloneqq\{B(v,1):v\in V_{G}\}$ of 1-neighborhood balls. Observe that a set $S$ is dominating in $G$ if and only if every vertex is covered by the subset of 1-neighborhood balls in $\mathcal{R}$ centered in $S$. In other words, a minimum solution of the set-cover instance $(V_{G},\mathcal{R})$ is a minimum dominating set in $G$. This simple observation immediately implies an $O(\log n)$-approximation to Dominating Set, where $n$ is the number of nodes in $G$, by the standard greedy set cover algorithm. $O(\log n)$-approximation is the best possible (under popular beliefs in complexity) if we make no further assumptions on the graph [25, 15]; this is essentially because the interaction between 1-neighborhood balls in general graphs can be complex and unruly.
\par
\par
% END S1.SS0.SSS0.P0.SPx1.p1
% BEGIN S1.SS0.SSS0.P0.SPx1.p2
What if we consider special classes of graphs with additional geometry constraints? One of the most well-studied examples is planar graphs, graphs with a drawing in the plane where no two edges cross each other. Planar graphs exhibit many helpful traits for algorithm design; particularly, several partitioning tools exist, including separators [43, 47, 59], low-diameter decompositions [41, 4, 9], and an algorithmic paradigm now famously known as Baker’s shifting/layering technique [1] which utilizes diameter-treewidth property [23] and the decomposability of planar graphs into $k$-outerplanar graphs. The idea is to compute a BFS-tree from an arbitrary node, then focus only on a few consecutive layers at a time. If the optimization problem at hand is local (for Dominating Set, node coverage can be checked by examining immediate neighbors), then a major portion of the optimal solution would be preserved within the consecutive layers. By iterating over the initial shift value and removing one out of every $O(1/\varepsilon)$ layers in the BFS-tree, one may guarantee a $(1+\varepsilon)$-approximation to the optimal solution. This is the main idea behind the polynomial-time approximation scheme (PTAS) of Dominating Set on unweighted planar graphs [1].
\par
\par
% END S1.SS0.SSS0.P0.SPx1.p2
% BEGIN S1.SS0.SSS0.P0.SPx1.p3
Not surprisingly, the same layering technique can be applied to the $r$-dominating set problem ($r$-Dominating Set for short) in planar graphs for any absolute constant $r$, where the goal is similar to Dominating Set but every node in $S$ can “cover” a vertex subset of radius $r$ [18]. $r$-Dominating Set was introduced by Barilan, Kortsarz, and Peleg [3] and later rediscovered in several other context [30, 18]. However, one encounters some trouble when generalizing Baker’s technique to minor-free graphs. This is because arbitrary minor-free graph no longer satisfies the diameter-treewidth property (imagine adding one extra vertex connecting to every node in a planar graph; the diameter becomes 2 but the graph is $K_{6}$-minor-free). One possible way to overcome the issue is to apply the Robertson-Seymour structural theorem for minor-free graphs [33, 58], which seems to be overkill. Another, possibly more satisfying solution, is to understand what the $O(1)$-radius neighborhood balls can look like in a minor-free graph. This is the route taken by the local search approach. On planar graphs, typical local search analysis [10] is usually done by applying an exchange argument on the base graph $G$, utilizing the fact that planar graphs have small separators [43] (in the form of an $r$-division [29]). Har-Peled and Quanrud [35] observed that, for problems like $r$-dominating set, the exchange argument in a typical local search analysis can be performed on the intersection graph of $O(1)$-neighborhood balls [48, 12] (more accurately, the subgraph $G_{\mathcal{R}}$ induced by the union of the local search solution and the optimal solution). They further observed that the intersection graph $G_{\mathcal{R}}$ over $O(1)$-radius balls in minor-free graph $G$, while not necessarily itself minor-free, must have polynomial expansion [50]. Without going into definitions, the class of graphs with polynomial expansion coincide with those that have sublinear-size separators for all subgraphs [21]. Thus one can apply the same kind of exchange argument using separator decomposition on the ball intersection graph $G_{\mathcal{R}}$.
\par
\par
% END S1.SS0.SSS0.P0.SPx1.p3
% END S1.SS0.SSS0.P0.SPx1
% BEGIN S1.SS0.SSS0.P0.SPx2
\par\medskip\noindent\textbf{Arbitrary radii and arbitrary weights.}\par
% BEGIN S1.SS0.SSS0.P0.SPx2.p1
The $r$-dominating set problem becomes significantly harder when $r$ can depend on $n$ [22]. In the literature this is usually refer to as the the metric/unbounded $r$-dominating set problem [26, 27], where we further allow arbitrary nonnegative edge-length $\ell$ on the input graph, in addition to the covering radius $r$ (now measuring with respect to $\ell$) becoming part of the input. We can even consider a weighted variant of the problem, where every vertex is associated with a positive-valued weight $w$, and the goal is to compute an $r$-dominating set with minimum total weight (instead of minimum cardinality). Baker’s layering technique no longer applies, nor does the local search strategy for the most part. (There is an exception [42, 26] which we will talk about in the technical section.) This is somewhat expected because the unbounded $r$-dominating set problem is no longer “local”, and the intersection graph of radius-$r$ balls may be dense and have complex interactions. In fact, under Exponential-Time Hypothesis, Fox-Epstein et al. [44, 28] proved that no efficient PTAS can exist for the unbounded $r$-dominating set problem on planar graphs. Thus, recent work has resorted to bicriteria algorithms, allowing the returned solution to use a dominating set of radius $(1+\varepsilon)r$ instead, in addition to the total weight being a $1+\varepsilon$ approximation to the optimal value [22, 28]. The idea is to embed the original planar/minor-free graph metric into some other tree-like metrics, where the $r$-dominating set problem can easily be solved using dynamic programming. Unfortunately one has to suffer some form of distortion (which translates to the slack in covering radius); the best result so far gives $+\varepsilon\Delta$-additive embedding of diameter-$\Delta$ minor-free graphs into graphs of $O_{\varepsilon}(1)$ treewidth [13, 14]. This line of research has been fruitful and ultimately leads to bicriteria EPTAS for planar graphs [28] and bicriteria Q(Q)PTAS for minor-free graphs [26, 27].
\par
\par
% END S1.SS0.SSS0.P0.SPx2.p1
% BEGIN S1.SS0.SSS0.P0.SPx2.p2
But what if we really want to solve the unbounded $r$-dominating set problem without compromise on the radius $r$? The lower bound result [44, 28] leaves us no choice but to look at $O(1)$-approximation (or at least, $(1+\varepsilon)$-approximation with running time $n^{f(\varepsilon)}$). The existence of such algorithm was explicitly asked by Filtser and Le [27].
\par
\par
% END S1.SS0.SSS0.P0.SPx2.p2
% END S1.SS0.SSS0.P0.SPx2
% BEGIN S1.SS0.SSS0.P0.SPx3
\par\medskip\noindent\textbf{Our contribution.}\par
% BEGIN S1.SS0.SSS0.P0.SPx3.p1
We establish the first polynomial-time $O(1)$-approximation algorithm for the (vertex)-weighted metric $r$-dominating set problem for unbounded $r$. We emphasize that $r$ is part of the input, and can be arbitrarily large compared to $n$.
\par
\par
% END S1.SS0.SSS0.P0.SPx3.p1
% BEGIN Thmtheorem1
\par\medskip\noindent\textbf{Theorem 1.}\par
Let $G$ be an arbitrary $n$-node planar graph, where each vertex $v$ is associated with a weight $w(v)$. Let $r$ be an arbitrary radius parameter. There is a polynomial-time algorithm that computes an $r$-dominating set in $G$ whose weight is $O(1)$ times the optimum.
\par
\par
\par
% END Thmtheorem1
% END S1.SS0.SSS0.P0.SPx3
% BEGIN S1.SS1
\par\medskip\noindent\textbf{1.1 Techniques}\par
% BEGIN S1.SS1.p1
Our new (single-criteria) $O(1)$-approximation algorithm for weighted $r$-dominating set on planar graphs uses the quasi-uniformity sampling technique, originally designed for weighted set cover problem on geometric shapes [63, 11]. But first, we have to describe a bit of (technical) history and two immediate challenges when adapting to our setting.
\par
\par
% END S1.SS1.p1
% BEGIN S1.SS1.SSS0.P0.SPx1
\par\medskip\noindent\textbf{Challenge 1: Planar support.}\par
% BEGIN S1.SS1.SSS0.P0.SPx1.p1
Geometric set cover problems have their own history and success [7, 16, 17, 62]. Notably, for an unweighted set cover instance $(X,\mathcal{R})$, the result of Brönnimann and Goodrich [7] provided a recipe to achieve $O(1)$-approximation. Their pioneering use of multiplicative weight update algorithm requires the existence of small $\varepsilon$-net: First, find an $\varepsilon$-net $N$ for the dual set system $(\mathcal{R},X)$ (notice that here the $\varepsilon$-net is in the dual and thus a subset of ranges in $\mathcal{R}$). If there are element $x$ in $X$ not covered by $N$, double the weight of every range in $\mathcal{R}$ that covers $x$, to increase the chance of choosing one such range in the next round. As a result, by setting $\varepsilon=\Theta(1/\mathrm{opt})$, this implies an $O(g(\mathrm{opt}))$-approximation for any set cover instance, given the dual system has an $\varepsilon$-net of size $O(\frac{1}{\varepsilon}\cdot g(\frac{1}{\varepsilon}))$. (See [24] for an LP interpretation and a simpler proof.)
\par
\par
% END S1.SS1.SSS0.P0.SPx1.p1
% BEGIN S1.SS1.SSS0.P0.SPx1.p2
For set system of bounded VC-dimension [61], there is always an $\varepsilon$-net of size $O(\frac{1}{\varepsilon}\log\frac{1}{\varepsilon})$ [36]. Also, if the VC-dimension of a primal system is $d$, then the dual system has VC-dimension at most $2^{d}$ [34]. To get $O(1)$-approximation, however, we need a much stronger assumption of $O(\frac{1}{\varepsilon})$-size $\varepsilon$-net in the dual system. This is not always achievable even for simple geometric shapes like lines and rectangles [49]; in fact, among the basic geometric shapes, only a handful are known to have linear-size $\varepsilon$-net, including halfspaces in 2D and 3D [46, 51, 45], unit-cubes in 3D [17], and family of pseudodisks in the plane (both primal and dual) [53].
\par
\par
% END S1.SS1.SSS0.P0.SPx1.p2
% BEGIN S1.SS1.SSS0.P0.SPx1.p3
For our application of unbounded $r$-dominating set problem, the “geometric shapes” are the radius-$r$ balls in planar graphs, which is discrete in nature, lacking well-defined boundary curves to even make the definition of pseudodisks meaningful. One might sneer at the issue and say, “there must be a way; how can distance balls in the plane be much different from pseudodisks?” But even if we redefine the notion of pseudodisks combinatorially (for every pair of regions $A$ and $B$, all three sets $A\cap B$, $A\setminus B$, and $B\setminus A$ are connected), the distance balls are still not pseudodisk. In fact, they can be piercing in the sense of Raman-Ray [54], as illustrated by Figure 1.
\par
\par
% END S1.SS1.SSS0.P0.SPx1.p3
% BEGIN S1.F1
\begin{center}\includegraphics[width=0.7\linewidth,height=0.38\textheight,keepaspectratio]{voronoi-page-8.pdf2svg.pdf}\end{center}
 Figure 1: Two distance balls $R$ and $R^{\prime}$ that are piercing. Here both $R$ and $R^{\prime}$ have radius 6. 
% END S1.F1
% BEGIN S1.SS1.SSS0.P0.SPx1.p4
Fortunately, Pyrga and Ray [53] provided an alternative method to obtain linear-size $\varepsilon$-net (in the dual) – to construct a sparse support for the primal system $(X,\mathcal{R})$. We say a set system $(X,\mathcal{R})$ has a support, if for any subsystem $(X^{\prime},\mathcal{R}_{|X^{\prime}})$ of $(X,\mathcal{R})$, there is an undirected graph $H$ on the elements of $X^{\prime}$ such that for every set $R$ in $\mathcal{R}_{|X^{\prime}}$, the subgraph $H_{R}$ of $H$ induced by the elements of $R$ is connected. The support is sparse if every such graph $H$ is sparse. The work of Pyrga and Ray [53] implies that if $(X,\mathcal{R})$ has bounded VC-dimension and a sparse support, then the dual system $(X,\mathcal{R})$ has $\varepsilon$-net of size $O(\frac{1}{\varepsilon})$.
\par
\par
% END S1.SS1.SSS0.P0.SPx1.p4
% BEGIN S1.SS1.SSS0.P0.SPx1.p5
Our first technical contribution is to compute a support for any collection of distance balls in any arbitrary graph (not just planar graphs). Crucially, if the given graph is planar, then so is the support. We remind the reader that our original goal is to solve the $r$-dominating set problem with unbounded $r$; thus the construction of the support has to hold for balls of unbounded radii as well. This is the first construction of its kind; most previous work on sparse support focus on geometric shapes, and almost entirely restricted to planar support for set systems that are pseudodisks [53] and their generalizations [54]. For piercing systems like axis-parallel rectangles, there are configurations where no planar support can exist [52].
\par
\par
% END S1.SS1.SSS0.P0.SPx1.p5
% BEGIN S1.SS1.SSS0.P0.SPx1.p6
Instead of focusing on planarity, we shift the narrative by arguing that a support graph for $(X,\mathcal{R})$ is really a connectivity sketch that exists for every graph $G$ on $X$, as long as the sets in $\mathcal{R}$ are distance balls on $G$. Our construction for support on arbitrary balls in $G$ can be described simply as follows: Build a Voronoi diagram of $V(G)$ with respect to balls centers, then contract every Voronoi cell into a single node. We argue that the contracted graph is a support of $G$. (The benefit of this additional property becomes evident when we tackle the next challenge.) Both planarity and sparseness come as a byproduct of the fact that contraction of the Voronoi diagram is a minor of $G$. (See Theorem 2 for a precise statement.) (A similar Voronoi diagram idea can be found in the work of Le [42, 26]; but their version only enables single-criteria algorithm for unweighted metric $r$-dominating set problem.)
\par
\par
% END S1.SS1.SSS0.P0.SPx1.p6
% END S1.SS1.SSS0.P0.SPx1
% BEGIN S1.SS1.SSS0.P0.SPx2
\par\medskip\noindent\textbf{Challenge 2: Linear shallow cell complexity.}\par
% BEGIN S1.SS1.SSS0.P0.SPx2.p1
There is another, perhaps more severe, obstacle to overcome. We want to solve the weighted version of the $r$-dominating set problem. The Brönnimann-Goodrich MWU approach relies on computing regular $\varepsilon$-net (as subset of balls covering the vertices); the size bound on $\varepsilon$-net translates to approximation guarantee on size of the cover, and thus can only work in the unweighted setting. Varadarajan [63] extended the $\varepsilon$-net idea to the weighted setting. This quasi-uniformity sampling technique computes $\varepsilon$-nets that samples each set in $\mathcal{R}$ with (roughly) equal probability, which then can be made aware of the weights on $\mathcal{R}$. Chan, Grant, Könemann, and Sharpe [11] later polished this technique and made connection to a more refined measure on set systems called shallow cell complexity. In our context, cell complexity refers to the number of different equivalent classes (called cells) of points that are covered by the same subset of $r$-balls. We say a cell has depth $k$ if the points within are covered by up to $k$ balls. A set system has shallow cell complexity $f(|\mathcal{R}|,k)$ if the number of cells of depth at most $k$ is bounded by $f(|\mathcal{R}|,k)$; for the purpose of obtaining $O(1)$-approximation for set cover, one requires the shallow cell complexity to be $O(|\mathcal{R}|\cdot\operatorname{poly}k)$.
\par
\par
% END S1.SS1.SSS0.P0.SPx2.p1
% BEGIN S1.SS1.SSS0.P0.SPx2.p2
The quasi-uniformity sampling technique can be used to solve the constant $r$ case for weighted $r$-dominating set problem on minor-free graphs (and graphs with bounded expansion as well [20]). This is because in the case of constant radius $r$, the whole cell complexity is linearly bounded [57, 38], which directly bounds the shallow counterpart. Unfortunately, it is known that the cell complexity of $r$-balls on a planar graph can be cubic in $r$ [38], and thus cannot help in the unbounded $r$ setting.
\par
\par
% END S1.SS1.SSS0.P0.SPx2.p2
% BEGIN S1.SS1.SSS0.P0.SPx2.p3
One way to see why halfspaces and pseudodisks are special is that they both have linear union complexity [40, 17]. Application of the classical Clarkson-Shor technique then implies that the shallow cell complexity is bounded by $O(|\mathcal{R}|\cdot k)$. When working with arbitrary-radius balls on graphs however, we run into trouble applying Clarkson-Shor. First, there is no well-defined notion of union complexity for a collection of subgraphs. One would imagine the support can take the role as the linear-size object at the base case for Clarkson-Shor; but this intuition is insufficient. The main issue is that for pseudodisks, the number of cells can be charged to the complexity of the whole pseudodisk arrangement by Euler’s formula, and every crossing in the arrangement is uniquely defined by two pseudodisks. There is no immediate equivalent object for distance balls that allows for such counting.
\par
\par
% END S1.SS1.SSS0.P0.SPx2.p3
% BEGIN S1.SS1.SSS0.P0.SPx2.p4
Our solution is to reimagine the unique representative argument by Raman and Ray for the non-piercing family [55]; the new proof is substantially different and relies on geometric properties of distance balls, instead of topological properties of curve sweeping [55, §6]. First we argue that every cell can be charged to a unique 3-tuple of ball centers (Lemma 7). This allows us to apply Clarkson-Shor and reduce the problem to counting cells of depth 1, 2, and 3. Depth 1 and 2 are straightforward; for depth-3 cells, we exploit the additional property that our support for distance balls came from contracting Voronoi cells with respect to the ball centers, and thus the number of depth-3 cells can be counted towards the complexity of the Voronoi diagram, which is asymptotically the same as the planar support (Lemma 13).
\par
\par
% END S1.SS1.SSS0.P0.SPx2.p4
% END S1.SS1.SSS0.P0.SPx2
% END S1.SS1
% END S1
% BEGIN S2
\par\medskip\noindent\textbf{2 Terminology and Support Graphs}\par
% BEGIN S2.p1
Let’s first introduce the terminology needed for concrete discussion. Our support construction actually works for balls in an arbitrary directed graph. We refer to the digraph that our ball system lives on as the underlying graph $G$. Note that all results apply to balls in undirected graphs. A ball $D$ can be defined using a pair $(c_{D},r_{D})$ where $c_{D}$ is a vertex of $G$ called the center, and a radius $r_{D}$ which is a real number. Ball $D$ is then defined to be the set of all vertices $v$ whose shortest path from $c_{D}$ has length at most $r_{D}$.
\par
\par
% END S2.p1
% BEGIN S2.p2
Given a system $(U,\mathcal{R})$, the dual system $(\mathcal{R},U)$ has the elements of $U$ treated as subsets of $\mathcal{R}$: if element $e$ in $U$ is a member of set $R$ in $\mathcal{R}$ in the primal system, then $R$ is a member of set $e$ in the dual system. We use the term dual support of $(U,\mathcal{R})$ to refer to a support graph of the dual system $(\mathcal{R},U)$. (We give an example of a dual support in Figure 2.) Although our $r$-dominating set application demands primal support, the set system of balls of equal radius is self-dual, Thus, we do not differentiate the support for the primal from the dual, but instead describe the support for the dual system since it is conceptually easier.
\par
\par
% END S2.p2
% BEGIN S2.F2
\begin{center}\includegraphics[width=0.7\linewidth,height=0.38\textheight,keepaspectratio]{voronoi-page-1.pdf2svg.pdf}\end{center}
 Figure 2: A set system with a dual support. The red regions containing the violet vertex $v$ form a connected subgraph in the support. 
% END S2.F2
% BEGIN S2.p3
In Section 3 we construct a dual support for any system of balls, by contracting cells of the Voronoi diagram of the underlying graph (with respect to ball centers) into a minor.
\par
\par
% END S2.p3
% BEGIN Thmtheorem2
\par\medskip\noindent\textbf{Theorem 2 (Support).}\par
Any system of balls $\mathcal{R}$ on an underlying digraph $G$ has a dual support $H(\mathcal{R})$ such that if $G$ does not contain a minor $K$ with minimum vertex degree two, then neither does $H(\mathcal{R})$.
\par
\par
\par
% END Thmtheorem2
% BEGIN Thmremarkx1
\par\medskip\noindent\textbf{$\blacktriangleright$ Remark.}\par
Unfortunately the min-degree requirement cannot be dropped in general. For example, if $G$ is $P_{2}$-minor-free, we have a set of disjoint vertices. But two balls can centered on the same vertex, and thus connected in the dual support, which implies the existence of an edge and so the support is not $P_{2}$-minor-free. Still the restriction is not severe as most of the interesting minor-free graph classes satisfy the min-degree condition, including $K_{h}$-minors for $h\geq 3$, grid minors, and all forbidden minors for bounded-genus graphs.
\par
\par
\par
% END Thmremarkx1
% BEGIN S2.p4
Since planar graphs can be classified by their excluded minors, these results imply a planar dual support for the ball system $\mathcal{R}$ on $G$. We then proceed to show that $\mathcal{R}$ has linear shallow cell complexity in Section 4.
\par
\par
% END S2.p4
% BEGIN Thmtheorem3
\par\medskip\noindent\textbf{Theorem 3 (Shallow Cell Complexity).}\par
Let $G$ be a planar graph. Any system of balls $(V_{G},\mathcal{R})$ has shallow cell complexity $O(|\mathcal{R}|\cdot k^{2})$.
\par
\par
\par
% END Thmtheorem3
% BEGIN S2.SS1
\par\medskip\noindent\textbf{2.1 Applications}\par
% BEGIN S2.SS1.p1
In addition to the result on computing unbounded $r$-dominating set in planar graphs, the supports also have their own implications. An important property of our support construction is that we only ever contract edges that are within radius $r$ of a ball center. This means that if $G$ is an unweighted graph with polynomial expansion and we are considering constant-radius balls, the support graph will also have polynomial expansion, allowing us to use the sublinear separators from Har-Peled and Quanrud [35] to get the same PTAS results for unweighted problems as Raman and Ray [54] achieved for non-piercing regions. If $r$ is not bounded by a constant, we can achieve the same when $G$ is minor-free [65].
\par
\par
% END S2.SS1.p1
% BEGIN S2.SS1.p2
Another natural use of supports is for problems relating to geodesic disks in a polygonal domain (possibly with holes). It is known that planar metrics and polygonal domains are equivalent: any system of geodesic disks in a polygonal domain in $\mathbb{R}^{2}$ can be represented as a system of balls in a planar graph, and vice versa [6]. By the equivalence, our techniques naturally give planar supports for geodesic disks in polygonal domains.
\par
\par
% END S2.SS1.p2
% END S2.SS1
% BEGIN S2.SS2
\par\medskip\noindent\textbf{2.2 Related Work on Support Graphs}\par
% BEGIN S2.SS2.p1
Van Cleemput [60] and Voloshina and Feinberg [64] used the existence of a planar support graph to give meaningful notion of planarity in the hypergraph setting. As such, support graphs have seen interest in (hyper)graph drawing community [37, 39, 8]. In the realm of optimization, Pyrga and Ray [53] used support graphs to construct linear-sized $\varepsilon$-nets for several set systems with bounded VC dimension, including halfspaces in $\mathbb{R}^{3}$ and pseudodisks (more generally, $r$-admissable regions) in $\mathbb{R}^{2}$. Here the support graphs were required to be sparse (i.e. with linearly many edges) in order to bootstrap existing $\varepsilon$-net sizes down to $O(\frac{1}{\varepsilon})$. Later on support graphs that are planar were used by Mustafa and Ray [48] to show a PTAS for various optimization problem in the geometric setting. Instead of $\varepsilon$-net, they exploited the existence of sublinear-size separators in planar graphs to show that an $O_{\varepsilon}(1)$-swap local search algorithm gives a $(1+\varepsilon)$-approximation for minimum geometric hitting set. The same paradigm of analyzing local search with separators on planar support was extended to other problems, such as maximum independent set and generalized set cover (including dominating set and hitting set) in pseudodisks and non-piercing regions [12, 19, 32, 54], terrain guarding [31, 2], and covering the boundary of a simple orthogonal polygon with rectangles [5]. (A survey on the many use of support can be found in Raman and Singh [56].)
\par
\par
% END S2.SS2.p1
% END S2.SS2
% END S2
% BEGIN S3
\par\medskip\noindent\textbf{3 Existence of Support Graph}\par
% BEGIN S3.p1
To begin constructing the dual support for a system of balls $\mathcal{R}$ in a graph $G$, we will first modify $G$ slightly and reduce to the case where every ball in $\mathcal{R}$ has the same radius. We will refer to the new graph as $G^{\prime}$. (Recall we work with the general case of digraphs.) We augment digraph $G$ by adding a new node corresponding to each ball in $\mathcal{R}$. (For clarity, we refer to original vertices of $G$ as vertices and new vertices corresponding to $\mathcal{R}$ as nodes.) For each ball $R$ in $\mathcal{R}$, we create a single node, $x_{R}$, which has only one edge: an outgoing edge into $c_{R}$, the center of the ball $R$. We give the edge a weight of $r_{\max}-r_{R}$, where $r_{\max}$ is defined to be $\max_{R^{\prime}\in\mathcal{R}}r_{R^{\prime}}$. This allows us to instead consider each $R$ in $\mathcal{R}$ to be centered at $x_{R}$, all with the same radius $r_{\max}$, and the balls will still define the same set system $\mathcal{R}$ over the vertices in $G$. Finally, remove all vertices which are not contained in any ball of $\mathcal{R}$ and slightly perturb the edge weights (without changing any ball containment relations) so that every vertex $v$ has a unique closest ball center in $\{x_{R}:R\in\mathcal{R}\}$.
\par
\par
% END S3.p1
% BEGIN S3.p2
To build the dual support, we construct a Voronoi partition of $G^{\prime}$ with respect to the set of nodes. A Voronoi partition $\mathcal{F}$ of a graph $G^{\prime}$ with respect to a set of nodes $S$ is a partition of $V_{G^{\prime}}$ into subsets such that for every node $R\in S$, there is a corresponding subset $\text{\emph{{\color[rgb]{.72,0,0}\definecolor[named]{pgfstrokecolor}{rgb}{%
.72,0,0}$F_{R}$}}}\in\mathcal{F}$ which contains all vertices $u$ in $G^{\prime}$ for which $d(x_{R},u)\leq\min_{x_{R}\in S\setminus\{x_{R}\}}d(x_{R^{\prime}},u)$. For clarity, we will refer to each set $F_{R}$ in $\mathcal{F}$ as the Voronoi cell of $R$. In our case, we construct the Voronoi partition $\mathcal{F}_{\mathcal{R}}$ of $G^{\prime}$ with respect to $\{x_{R}:R\in\mathcal{R}\}$, which we identify as $\mathcal{R}$. Each cell $F_{R}$ in the Voronoi partition is a connected subset of $G^{\prime}$ because for every vertex $v$ in $F_{R}$ reachable from $x_{R}$, the vertices on the shortest path between $x_{R}$ and $v$ must belong to $F_{R}$ as well. Thus we can simply contract every $F$ in $\mathcal{F}_{\mathcal{R}}$ to a (contraction)-minor on only the nodes in $\mathcal{R}$. This final graph $H$, we claim, is a dual support for $(G,\mathcal{R})$.
\par
\par
% END S3.p2
% BEGIN S3.F3
\begin{center}\includegraphics[width=0.7\linewidth,height=0.38\textheight,keepaspectratio]{voronoi-page-3.pdf2svg.pdf}\end{center}
 Figure 3: An illustration of the proof to Lemma 4 that $\pi$ can only pass through Voronoi cells of $\mathcal{R}_{v}$. The purple region is the Voronoi cell $F_{Q}$. 
% END S3.F3
% BEGIN Thmtheorem4
\par\medskip\noindent\textbf{Lemma 4.}\par
For any vertex $v$ in $G^{\prime}$, the set of balls $\emph{{\color[rgb]{.72,0,0}\definecolor[named]{pgfstrokecolor}{rgb}{.72,0,0}$%
\mathcal{R}_{v}$}}\subseteq\mathcal{R}$ that contain $v$ must induce a connected subgraph in $H$.
\par
\par
\par
% END Thmtheorem4
% BEGIN S3.SS0.SSS0.P0.SPx1
\par\medskip\noindent\textbf{Proof.}\par
% BEGIN S3.SS0.SSS0.P0.SPx1.p1
Consider any vertex $v$ which is contained in at least two balls in $\mathcal{R}$. Vertex $v$ must be in the Voronoi cell $F_{R}$ of some ball $R$ by construction of $G^{\prime}$. We claim that for any other ball $R^{\prime}$ of $\mathcal{R}$ that also contains $v$, there is a path from $x_{R^{\prime}}$ to $v$ in $G^{\prime}$, denoted as $\pi$, that only passes through Voronoi cells corresponding to balls that contain $v$.
\par
\par
% END S3.SS0.SSS0.P0.SPx1.p1
% BEGIN S3.SS0.SSS0.P0.SPx1.p2
Since we have removed all vertices which have no incoming path from any ball, every vertex along $\pi$ in $G^{\prime}$ must be in some cell. Assume that some vertex $u$ along $\pi$ is contained in some Voronoi cell $F_{Q}$ corresponding to a ball $Q$ that does not contain $v$. By definition of Voronoi cell, we have that $d(x_{Q},u)\leq d(x_{R^{\prime}},u)$. By the definition of shortest paths, we know that $d(x_{R^{\prime}},v)=d(x_{R^{\prime}},u)+d(u,v)$ . By triangle inequality, we then have the following:
\par
\[\begin{aligned}
d(x_{Q},v)\leq d(x_{Q},u)+d(u,v)\leq d(x_{R^{\prime}},u)+d(u,v)=d(x_{R^{\prime%
}},v).
\end{aligned}\]
But, $d(x_{R^{\prime}},v)\leq r_{\max}$ and $d(x_{Q},v)>r_{\max}$ since $v$ is in $R^{\prime}$ but not in $Q$, contradiction. $\hfill\blacktriangleleft$
\par
\par
% END S3.SS0.SSS0.P0.SPx1.p2
% END S3.SS0.SSS0.P0.SPx1
% BEGIN S3.SS0.SSS0.P0.SPx2
\par\medskip\noindent\textbf{Proof of Theorem 2.}\par
% BEGIN S3.SS0.SSS0.P0.SPx2.p1
Construct $H$ as described above. By Lemma 4, we know that $H$ satisfies the definition of a dual support for system $(V_{G},\mathcal{R})$. It remains to show that if $G$ has no minor $K$ with minimum degree two, then neither can $H$. This is done by arguing that $G^{\prime}$ has no such minors, and then the fact that $H$ is a minor of $G^{\prime}$ concludes the proof.
\par
\par
% END S3.SS0.SSS0.P0.SPx2.p1
% BEGIN S3.SS0.SSS0.P0.SPx2.p2
When constructing $G^{\prime}$, we only perform two types of operations that would affect the existence of minors: deleting vertices and adding in nodes $\{x_{R}\}$ of degree one. Let $K$ be a minor model (as a collection of disjoint vertex subsets called supernodes) of $G^{\prime}$ that has minimum degree two. If $K$ only contains vertices but not nodes in $G^{\prime}$, then $K$ must appear in $G$ as well (even though $G^{\prime}$ was obtained by removing some vertices from $G$). So at least one supernode $\eta$ of $K$ must contain some degree-1 $x_{R}$ that was added in $G^{\prime}$.
\par
\par
% END S3.SS0.SSS0.P0.SPx2.p2
% BEGIN S3.SS0.SSS0.P0.SPx2.p3
For an edge $(\eta,\eta^{\prime})$ to exist in $K$, there must be a vertex/node $u\in\eta$ and a vertex/node $u^{\prime}\in\eta^{\prime}$ such that $(u,u^{\prime})$ is an edge of $G^{\prime}$. For $\eta$ containing $x_{R}$, $x_{R}$ must not be the only vertex/node in $\eta$, since $\eta$ have degree at least 2. Because $\eta$ is connected, the neighbor of $x_{R}$ is also in $\eta$ and the (sole) edge incident on $x_{R}$ must be contracted. This, however, means that there can be no edge $(x_{R},u^{\prime})$ in $G^{\prime}$ that corresponds to an edge $(\eta,\eta^{\prime})$ in $K$, so there is a minor $K^{\prime}$ isomorphic to $K$ where $x_{R}$ was simply deleted and thus $x_{R}$ is not in $\eta$. Since this holds for all $x_{R}$, there must be some minor of $G$ isomorphic to $K$, a contradiction. $\hfill\blacktriangleleft$
\par
\par
% END S3.SS0.SSS0.P0.SPx2.p3
% END S3.SS0.SSS0.P0.SPx2
% END S3
% BEGIN S4
\par\medskip\noindent\textbf{4 Shallow Cell Complexity of Balls in Planar Graphs}\par
% BEGIN S4.p1
Shallow cell complexity was introduced by Chan et al. [11] as a combinatorial analog of union complexity. We introduce the terminology within our context of distance balls on graphs.
\par
\par
% END S4.p1
% BEGIN S4.p2
Define the hit set of a vertex $v$ with respect to $\mathcal{R}$ to be the subset of balls in $\mathcal{R}$ which contain $v$. A cell of the system $(V_{G},\mathcal{R})$ is a set of vertices in $V_{G}$ that have the same hit set of incident balls. Cell complexity is simply the number of cells in $(V_{G},\mathcal{R})$, where as shallow cell complexity counts the number of cells of a particular depth $k$ in $\mathcal{R}$, denoted $\mathit{Cell}_{k}$, the equivalence classes of vertices whose hit set has size $k$. We say the set system has shallow cell complexity $f(n,k)$ if for depth $k$, there are at most $f(n,k)$ cells. Our goal for the section is to show that for a planar graph $G$, the shallow cell complexity for any system of balls $(V_{G},\mathcal{R})$ is $O(|\mathcal{R}|\cdot k^{2})$ (Theorem 3). Just as in the support constructions in § 3, we again reduce to the case where every ball in our system $\mathcal{R}$ has the same radius.
\par
\par
% END S4.p2
% BEGIN S4.SS0.SSS0.P0.SPx1
\par\medskip\noindent\textbf{Proof of Theorem 3.}\par
% BEGIN S4.SS0.SSS0.P0.SPx1.p1
We use two key lemmas proven later in the subsections. The first claims that for every cell $c$ of depth $k$, there is a unique encoding of $c$, $\langle\alpha(c),\beta(c),\gamma(c)\rangle$, where each member of the tuple is a ball of $\mathcal{R}$ (Lemma 7 in Section 4.1). Importantly, the encoding is unique among cells of depth $k$, not universally. This naively implies that there is at most $O(|\mathcal{R}|^{3})$ cells of depth $k$. Using the standard Clarkson–Shor arguments argument, we can improve this to $O(|\mathcal{R}|\cdot k^{2})$ if the number of cells of depth (at most) $3$ is $O(|\mathcal{R}|)$ (Lemma 13 in Section 4.2). More precisely: Sample the regions of $\mathcal{R}$ with probability $1/k$. A cell $c$ with unique encoding appears in the sample as depth-3 cell with probability at least $(ek)^{-3}$. The expected number of depth-3 cells in the sample is at most $O(|\mathcal{R}|/k)$, so the number of depth-$k$ cells in $\mathcal{R}$ is at most $O(|\mathcal{R}|\cdot k^{2})$. $\hfill\blacktriangleleft$
\par
\par
% END S4.SS0.SSS0.P0.SPx1.p1
% END S4.SS0.SSS0.P0.SPx1
% BEGIN S4.SS1
\par\medskip\noindent\textbf{4.1 Unique 3-Site Encoding of Cells}\par
% BEGIN S4.SS1.p1
We begin by assigning each cell of $\mathit{Cell}_{k}$ an unique encoding for each fixed depth $k\geq 3$. This encoding will depend on the set of incident balls corresponding to $c$. Specifically, the encoding will be formed by a 3-tuple of regions in $\mathcal{R}$, which we will refer to as $\langle\alpha,\beta,\gamma\rangle$. In assigning the encoding, we do not modify to the set system at all. Instead, we fix an arbitrary vertex $v_{c}$ for cell $c$ as its representative. We denote the cell containing $v$ as $\mathit{Cell}(v)$.
\par
\par
% END S4.SS1.p1
% BEGIN S4.SS1.p2
For a cell $c$, we begin by designating its $\langle\alpha,\beta\rangle$-type. We define $\alpha(v)$ as the furthest ball from $v$ that contains $v$. Similarly, we define $\beta(v)$ as the second furthest ball from $v$ that contains $v$. For cell $c$, its $\langle\alpha,\beta\rangle$-type is simply $\langle\alpha(v_{c}),\beta(v_{c})\rangle$. (Notice that $\alpha(v_{c})$ and $\beta(v_{c})$ might be different if one chooses a different $v_{c}$ from cell $c$, and thus the $\langle\alpha,\beta\rangle$-type of $c$ depends on its representative.) Let $\pi_{\alpha,\beta}(v)$ be the union of $\pi(\alpha(v),v)$ and $\pi(\beta(v),v)$, where $\pi(x,y)$ is the shortest path from $x$ to $y$ in $G$.
\par
\par
% END S4.SS1.p2
% BEGIN S4.F4
\begin{center}\includegraphics[width=0.7\linewidth,height=0.38\textheight,keepaspectratio]{voronoi-page-9.pdf2svg.pdf}\end{center}
 Figure 4: Given two vertices $v,v^{\prime}$ satisfying $\alpha(v)=\alpha(v^{\prime})=D_{1}$, $\beta(v)=\beta(v^{\prime})=D_{2}$ and $v^{\prime}\in\pi(D_{1},v)$, any ball $D_{3}$ containing $v^{\prime}$ must also contain $v$. 
% END S4.F4
% BEGIN Thmtheorem5
\par\medskip\noindent\textbf{Lemma 5.}\par
For a vertex $v$ with $\alpha(v)=D_{1}$ and $\beta(v)=D_{2}$, every vertex $v^{\prime}$ in $\pi_{\alpha,\beta}(v)$ with $\alpha(v^{\prime})=D_{1}$ and $\beta(v^{\prime})=D_{2}$ satisfies $\mathit{hit}(v^{\prime})\subseteq\mathit{hit}(v)$.
\par
\par
\par
% END Thmtheorem5
% BEGIN S4.SS1.SSS0.P0.SPx1
\par\medskip\noindent\textbf{Proof.}\par
% BEGIN S4.SS1.SSS0.P0.SPx1.p1
Assume the contrary. Then, there is some vertex $v^{\prime}$ which is contained within some ball $D_{3}$ such that $v\not\in D_{3}$. By definition of $\alpha$ and $\beta$, it must be that the center of $D_{3}$ is closer to $v^{\prime}$ than $D_{1}$ and $D_{2}$. However, since $v^{\prime}$ is either in $\pi(c_{D_{1}},v)$ or $\pi(c_{D_{2}},v)$, this means that $D_{3}$ is closer to $v$ than one of $D_{1}$ or $D_{2}$, and thus $D_{3}$ must contain $v$, a contradiction by the assumption that all balls have equal radius. (See Figure 4.) $\hfill\blacktriangleleft$
\par
\par
% END S4.SS1.SSS0.P0.SPx1.p1
% BEGIN S4.SS1.SSS0.P0.SPx1.p2
The third component of our encoding, $\gamma$, can now be defined. For $k=3$, $\gamma$ is simply the remaining ball after removing $\alpha$ and $\beta$. For $k\geq 4$, we define $\gamma$ procedurally after fixing $\alpha,\beta$. Denote the set of all cells of depth $k$ with $\alpha(v_{c})=D_{1}$ and $\beta(v_{c})=D_{2}$ as $C_{k}(D_{1},D_{2})$. We additionally collect the representatives $v_{c}$ of each $c\in C_{k}(D_{1},D_{2})$ into a corresponding set $V_{k}(D_{1},D_{2})$. Let $T_{D}(S)$ be the shortest path tree from the center $x_{D}$ of a ball $D$ to a set of vertices $S$. Consider the shortest path tree $T_{D_{1}}(V_{k}(D_{1},D_{2}))$. We will process all $v_{c}\in V_{k}(D_{1},D_{2})$ using the plane drawing of $T_{D_{1}}(V_{k}(D_{1},D_{2}))$.
\par
\par
% END S4.SS1.SSS0.P0.SPx1.p2
% BEGIN S4.SS1.SSS0.P0.SPx1.p3
Perform a depth-first traversal starting at $x_{D}$, processing children in clockwise order. The depth-first traversal $T_{D_{1}}(V_{k}(D_{1},D_{2}))$ of gives a total ordering of the vertices in $V_{k}(D_{1},D_{2})$, which also induces a total ordering $\sigma$ on the corresponding cells $C_{k}(D_{1},D_{2})$. For brevity, we will index the cells based on their order in $\sigma$ (the $i$th cell as $c_{i}$) in a cyclic manner. For the $i$th cell $c_{i}$ in the total ordering, we assign $\gamma(c_{i})$ to be an arbitrary ball $D$ with $D\in\mathit{hit}(v_{c_{i}})$ and $D\not\in\mathit{hit}(v_{c_{i-1}})$. Since $|\mathit{hit}(v_{c_{i}})|=|\mathit{hit}(v_{c_{i}-1})|=k$ and $\mathit{hit}(v_{c_{i}})\neq\mathit{hit}(v_{c_{i}-1})$ (by the definition of cell), such a $D$ must always exist, and thus we can assign $\gamma(c_{i})$ to all cells.
\par
\par
% END S4.SS1.SSS0.P0.SPx1.p3
% BEGIN S4.SS1.SSS0.P0.SPx1.p4
We now show that this definition of $\gamma$ is sufficient to ensure that each distinct cell $c$ of size $k$ has a unique encoding. First, we use Lemma 5 to show that the representative vertices of cells with the same $\alpha$ and $\beta$ cannot form ancestor relationships in $T_{D_{1}}(V_{k}(D_{1},D_{2}))$.
\par
\par
% END S4.SS1.SSS0.P0.SPx1.p4
% BEGIN Thmtheorem6
\par\medskip\noindent\textbf{Observation 6.}\par
For any two distinct cells $c,c^{\prime}\in C_{k}(D_{1},D_{2})$, $v_{c}$ cannot be an ancestor of $v_{c^{\prime}}$ in $T_{D_{1}}(V_{k}(D_{1},D_{2})$ or $T_{D_{2}}(V_{k}(D_{1},D_{2})$.
\par
\par
\par
% END Thmtheorem6
% END S4.SS1.SSS0.P0.SPx1
% BEGIN S4.SS1.SSS0.P0.SPx2
\par\medskip\noindent\textbf{Proof.}\par
% BEGIN S4.SS1.SSS0.P0.SPx2.p1
First consider $T_{D_{1}}(V_{k}(D_{1},D_{2})$. If $v_{c}$ were an ancestor of $v_{c^{\prime}}$, then $v_{c}$ would be in $\pi(D_{1},v_{c^{\prime}})$. This would imply that $\mathit{hit}(v_{c^{\prime}})\subseteq\mathit{hit}(v_{c})$ by Lemma 5. However, since $c$ and $c^{\prime}$ have depth $k$, $|\mathit{hit}(v_{c})|=|\mathit{hit}(v_{c^{\prime}})|=k$. The subset relationship would hence require that $\mathit{hit}(v_{c})=\mathit{hit}(v_{c^{\prime}})$, and thus $c$ and $c^{\prime}$ are the same class, which is a contradiction. A symmetric argument applies for $T_{D_{2}}(V_{k}(D_{1},D_{2})$. $\hfill\blacktriangleleft$
\par
\par
% END S4.SS1.SSS0.P0.SPx2.p1
% BEGIN S4.SS1.SSS0.P0.SPx2.p2
This means that for any $c,c^{\prime}$ with the same $\alpha$ and $\beta$, $\pi(\alpha(v_{c}),v_{c})\not\subseteq\pi(\alpha(v_{c}),v_{c^{\prime}})$.
\par
\par
% END S4.SS1.SSS0.P0.SPx2.p2
% BEGIN Thmtheorem7
\par\medskip\noindent\textbf{Lemma 7.}\par
For $k\geq 3$, every cell $c$ of depth $k$ has a unique encoding $\langle\alpha(c),\beta(c),\gamma(c)\rangle\in\mathcal{R}^{3}$.
\par
\par
\par
% END Thmtheorem7
% END S4.SS1.SSS0.P0.SPx2
% BEGIN S4.SS1.SSS0.P0.SPx3
\par\medskip\noindent\textbf{Proof.}\par
% BEGIN S4.SS1.SSS0.P0.SPx3.p1
If two cells have the same depth and $\alpha$ and $\beta$, they are both indexed by the same traversal order $\sigma$. Consider two cells $c_{i},c_{j}$ such that $\langle\alpha(c_{i}),\beta(c_{i}),\gamma(c_{i})\rangle=\langle\alpha(c_{j}),%
\beta(c_{j}),\gamma(c_{j})\rangle=(D_{1},D_{2},D_{3})$ and $|\mathit{hit}(v_{c_{i}})|=|\mathit{hit}(v_{c_{j}})|=k$.
\par
\par
% END S4.SS1.SSS0.P0.SPx3.p1
% BEGIN S4.SS1.SSS0.P0.SPx3.p2
We need to consider how $c_{i}$ and $c_{j}$ were assigned the same $\gamma$. Recall that we force adjacent cells $c_{i-1},c_{i}$ to not just have different $\gamma$ assignments, but that the $\gamma(c_{i})$ not contain $\gamma(c_{i-1})$. This forces $v_{c_{i-1}}$ and $v_{c_{j-1}}$ to not be contained by $D_{3}$. Let $a$ be the lowest common ancestor of $v_{c_{i-1}}$ and $v_{c_{j-1}}$ in $T_{D_{1}}(V_{k}(D_{1},D_{2})$. Similarly, let $b$ be the lowest common ancestor of $v_{c_{i-1}}$ and $v_{c_{j-1}}$ in $T_{D_{2}}(V_{k}(D_{1},D_{2})$. We consider the cycle $F$ formed by the union of the following shortest paths: $\pi(a,v_{c_{i-1}})$ to $\pi(v_{c_{i-1}},b)$ to $\pi(b,v_{c_{j-1}})$ to $\pi(v_{c_{j-1}},a)$. Since $G$ is planar, $F$ describes the boundary of some region $R(F)$. Since $i-1<i<j-1<j$, by the definition of our traversal one of $v_{c_{i}}$ or $v_{c_{j}}$ lies in $R(F)$ and the other lies outside. (We give an example $R(F)$ as a gray region in Figure 5.) This means that, for $D_{3}$ to contain both $v_{c_{i}}$ and $v_{c_{j}}$, $D_{3}$ must also contain some vertex $u^{\prime}$ of $F$. Assume WLOG that $u^{\prime}$ lies on $\pi(D_{3},v_{c_{j}})$. Since $u^{\prime}$ lies on $F$, it must be on one of $\pi(D_{1},v_{c_{i-1}}),\pi(D_{1},v_{c_{j-1}}),\pi(D_{2},v_{c_{i-1}}),$ or $\pi(D_{2},v_{c_{i-1}})$.
\par
\par
% END S4.SS1.SSS0.P0.SPx3.p2
% BEGIN S4.SS1.SSS0.P0.SPx3.p3
\begin{enumerate}
\item[]  $\blacksquare$ 
If $u^{\prime}$ is in $\pi(D_{1},v_{c_{i-1}})$ or $\pi(D_{1},v_{c_{j-1}})$, then we know that $D_{3}$ must be further from $u^{\prime}$ than $D_{1}$ since $D_{3}$ cannot contain $v_{c_{i-1}}$ or $v_{c_{j-1}}$. However, since $\alpha(v_{c_{j}})=D_{1}$, $D_{3}$ must be closer to $u^{\prime}$ than $D_{1}$ since $u^{\prime}$ is on $\pi(D_{3},v_{c_{j}})$ and $\alpha(v_{c_{j}})$ is the furthest ball from $v_{c_{j}}$ which still contains it. (See left subfigure of Figure 5.)
\par
\par
\item[]  $\blacksquare$ 
Similarly, if $u^{\prime}$ is in $\pi(D_{2},v_{c_{i-1}})$ or $\pi(D_{2},v_{c_{j-1}})$, then $D_{3}$ must be further from $u^{\prime}$ than $D_{2}$ since $D_{3}$ cannot contain $v_{c_{i-1}}$ or $v_{c_{j-1}}$. However, since $\beta(v_{c_{j}})=D_{2}$, $D_{3}$ must be closer to $u^{\prime}$ than $D_{2}$ since $u^{\prime}$ is on $\pi(D_{3},v_{c_{j}})$ and $\beta(v_{c_{j}})$ is the second furthest ball from $v_{c_{j}}$, while $D_{3}$ is at best third furthest. (See right subfigure of Figure 5.)
\par
\par
\end{enumerate}
\par
% END S4.SS1.SSS0.P0.SPx3.p3
% BEGIN S4.F5
\begin{center}\includegraphics[width=0.7\linewidth,height=0.38\textheight,keepaspectratio]{voronoi-page-10.pdf2svg.pdf}\end{center}
\begin{center}\includegraphics[width=0.7\linewidth,height=0.38\textheight,keepaspectratio]{voronoi-page-11.pdf2svg.pdf}\end{center}
  Figure 5: If two cells $c_{i},c_{j}$ have the same $\gamma(c_{i})=\gamma(c_{j})=D_{3}$, then there must $D_{3}$ must contain a vertex $u^{\prime}$ of $\gamma_{\alpha,\beta}(v_{c_{j-1}})$ or $\gamma_{\alpha,\beta}(v_{c_{j-1}})$. The region $R(F)$ corresponding to the simple cycle $F$ is given in gray. Left: $u^{\prime}$ is in $\pi(D_{1},v_{c_{j-1}})$. Right: $u^{\prime}$ is in $\pi(D_{2},v_{c_{j-1}})$. 
% END S4.F5
% BEGIN S4.SS1.SSS0.P0.SPx3.p4
We get a contradiction in both cases, which means that $\gamma(c_{i})$ cannot equal $\gamma(c_{j})$, and thus every cell $c$ of the same depth $k\geq 3$ must have a unique encoding $\langle\alpha(c),\beta(c),\gamma(c)\rangle$. $\hfill\blacktriangleleft$
\par
\par
% END S4.SS1.SSS0.P0.SPx3.p4
% END S4.SS1.SSS0.P0.SPx3
% END S4.SS1
% BEGIN S4.SS2
\par\medskip\noindent\textbf{4.2 Depth-3 Cells}\par
% BEGIN S4.SS2.p1
Now that we have a unique encoding $\langle\alpha(v_{c}),\beta(v_{c}),\gamma(v_{c})\rangle$ for every cell $c$ of depth $k$, we can proceed to give a bound on $|\mathit{Cell}_{3}|$.
\par
\par
% END S4.SS2.p1
% BEGIN S4.SS2.p2
For each $c$ in $\mathit{Cell}_{3}$, we consider its encoding $\langle\alpha(v_{c}),\beta(v_{c}),\gamma(v_{c})\rangle$. Since $c$ only has depth 3, we know that these are simply the balls of $\mathit{hit}(v_{c})$, placed in order of distance from $v_{c}$. Our first step is to select one of the three balls of $\mathit{hit}(v_{c})$ as $\mathit{base}(v_{c})$. Next, we iterate through every ball $D$ in $\mathcal{R}$ and count the number of cells that have been assigned $D$. We accomplish this by constructing a planar graph $G_{D}$ on the neighbors of $D$ in the Voronoi partition graph $H$. Summing over all $D$, this yields a total count of $O(|\mathcal{R}|)$ cells. We start with an important lemma regarding the order of distances to balls from vertices that lie on a shortest path.
\par
\par
% END S4.SS2.p2
% BEGIN Thmtheorem8
\par\medskip\noindent\textbf{Lemma 8.}\par
Given any ball $D$ of $\mathcal{R}$ , a vertex $v$, some vertex $v^{\prime}$ on $\pi(x_{D},v)$, and a distinct ball $D^{\prime}\in\mathcal{R}$, if $d(x_{D^{\prime}},v^{\prime})\leq d(x_{D},v^{\prime})$, then for all $v^{\prime\prime}\in\pi(v^{\prime},v)$, $d(x_{D^{\prime}},v^{\prime\prime})\leq d(x_{D},v^{\prime\prime})$.
\par
\par
\par
% END Thmtheorem8
% BEGIN S4.SS2.SSS0.P0.SPx1
\par\medskip\noindent\textbf{Proof.}\par
% BEGIN S4.SS2.SSS0.P0.SPx1.p1
For any $v^{\prime\prime}\in\pi(v^{\prime},v)$, because $v^{\prime}$ is between $x_{D}$ and $v^{\prime\prime}$ on $\pi(x_{D},v)$, $d(x_{D},v^{\prime\prime})=d(x_{D},v^{\prime})+d(v^{\prime},v^{\prime\prime})$. We also know that $d(x_{D^{\prime}},v^{\prime\prime})\leq d(x_{D^{\prime}},v^{\prime})+d(v^{%
\prime},v^{\prime\prime})$. Since $d(x_{D^{\prime}},v^{\prime})\leq d(x_{D},v^{\prime})$, we can simply plug in to get $d(x_{D^{\prime}},v^{\prime\prime})\leq d(x_{D},v^{\prime})+d(v^{\prime},v^{%
\prime\prime})=d(x_{D},v^{\prime\prime})$. $\hfill\blacktriangleleft$
\par
\par
% END S4.SS2.SSS0.P0.SPx1.p1
% BEGIN Thmtheorem9
\par\medskip\noindent\textbf{Observation 9.}\par
For a cell $c$ and ball $D^{\prime}\in\mathit{hit}(v_{c})$, all vertices $v$ in $\pi(x_{D},v_{c})$ must lie in the Voronoi cell of a ball in $\mathit{hit}(v_{c})$.
\par
\par
\par
% END Thmtheorem9
% END S4.SS2.SSS0.P0.SPx1
% BEGIN S4.SS2.SSS0.P0.SPx2
\par\medskip\noindent\textbf{Proof.}\par
% BEGIN S4.SS2.SSS0.P0.SPx2.p1
Consider a ball $D^{\prime}\not\in\mathit{hit}(v_{c})$ for which a vertex $v$ of $\pi(x_{D},v_{c})$ lies in the Voronoi cell $F_{D}$. If this is the case, then all the vertices after $v$ in $\pi(x_{D},v_{c})$ will be closer to $D^{\prime}$ than they are to $D$, and since $D$ contains $v_{c}$, $D^{\prime}$ must also then contain $v_{c}$, which is a contradiction. $\hfill\blacktriangleleft$
\par
\par
% END S4.SS2.SSS0.P0.SPx2.p1
% BEGIN S4.SS2.SSS0.P0.SPx2.p2
This means that the shortest paths from the centers of the balls hit by $v_{c}$ to the vertex $v_{c}$ must stay within the Voronoi cells of $\mathit{hit}(v_{c})$. This observation, paired with Lemma 8, allows us to define $\mathit{base}(v_{c})$ based on $\pi(x_{\alpha(v_{c})},v_{c})$. We consider two cases:
\par
\begin{enumerate}
\item[]  $\blacksquare$ 
Case 1: Some vertex $v$ in $\pi(x_{\alpha(v_{c})},v_{c})$ lies in $F_{\beta(v_{c})}$.  Assign $\mathit{base}(v_{c})\leftarrow\beta(v_{c})$.
\par
\par
\item[]  $\blacksquare$ 
Case 2: No vertex $v$ in $\pi(x_{\alpha(v_{c})},v_{c})$ lies in $F_{\beta(v_{c})}$.  Assign $\mathit{base}(v_{c})\leftarrow\gamma(v_{c})$.
\par
\par
\end{enumerate}
\par
% END S4.SS2.SSS0.P0.SPx2.p2
% BEGIN S4.F6
\begin{center}\includegraphics[width=0.7\linewidth,height=0.38\textheight,keepaspectratio]{voronoi-page-19.pdf2svg.pdf}\end{center}
\begin{center}\includegraphics[width=0.7\linewidth,height=0.38\textheight,keepaspectratio]{voronoi-page-18.pdf2svg.pdf}\end{center}
  Figure 6: Showing the subwalk of Lemma 10 and how to select $\mathit{base}$. Left: Case 1. Right: Case 2. 
% END S4.F6
% BEGIN S4.SS2.SSS0.P0.SPx2.p3
Now that we have assigned every cell to a ball, we consider a ball $D$ and the set $\mathit{resident}(D)$, which is the set of cells $c$ with $\mathit{base}(c)=D$. We will consider the cells of $\mathit{resident}(D)$ in relation to the neighbors of $D$ within the Voronoi partition graph $H$.
\par
\par
% END S4.SS2.SSS0.P0.SPx2.p3
% BEGIN Thmtheorem10
\par\medskip\noindent\textbf{Lemma 10.}\par
For a ball $D$ and every cell $c$ of $\mathit{resident}(D)$, there is a subwalk $\pi_{c}$ of $\pi_{\alpha,\beta}(v_{c})$ with the (nonempty set of) interior vertices lying in $D$ and the end vertices lying in distinct neighbors $D^{\prime},D^{\prime\prime}$ of $D$ in $H$.
\par
\par
\par
% END Thmtheorem10
% END S4.SS2.SSS0.P0.SPx2
% BEGIN S4.SS2.SSS0.P0.SPx3
\par\medskip\noindent\textbf{Proof.}\par
% BEGIN S4.SS2.SSS0.P0.SPx3.p1
If cell $c$ is a case-1 cell, we have that a vertex $v$ in $\pi(x_{\alpha(v_{c})},v_{c})$ which lies in $F_{\beta(v_{c})}$. By Lemma 8, no vertex of $\pi(v,v_{c})$ can lie in $F_{\alpha(v_{c})}$. Additionally, $v_{c}$ must lie in $F_{\gamma(v_{c})}$ by definition of $\alpha$ and $\beta$. Let $u$ be the first vertex of $\pi(v,v_{c})$ which lies in $F_{\gamma(v_{c})}$. Let $v^{\prime}$ be last vertex of $\pi_{\alpha,\beta}(v_{c})$ which lies in $F_{\alpha(v_{c})}$. The subwalk $\pi_{\alpha,\beta}(v_{c})[v^{\prime},u]$ satisfies this condition.
\par
\par
% END S4.SS2.SSS0.P0.SPx3.p1
% BEGIN S4.SS2.SSS0.P0.SPx3.p2
If cell $c$ is a case-2 cell, we have no such vertex and $\pi(x_{\alpha(v_{c})},v_{c})$ lies entirely in $\alpha(v_{c})$ or $\gamma(v_{c})$. Let $u$ be the last vertex of $\pi(x_{\alpha(v_{c})},v_{c})$ which lies in $F_{\alpha(v_{c})}$ and let $u^{\prime}$ be the last vertex of $\pi(x_{\beta(v_{c})},v_{c})$ that lies in $F_{\beta(v_{c})}$. We then can take our subpath to be $\pi_{\alpha,\beta}(v_{c})[u,u^{\prime}]$, and this will satisfy the condition. $\hfill\blacktriangleleft$
\par
\par
% END S4.SS2.SSS0.P0.SPx3.p2
% BEGIN S4.SS2.SSS0.P0.SPx3.p3
We now consider the neighbors of $D$ in the Voronoi support $H$, $N_{H}(D)$. In Lemma 10, the subwalk $\pi_{c}$ of cell $c$ in $\mathit{resident}(D)$ runs between the two balls of $\mathit{hit}(v_{c})\setminus D$, both of which must be in $N_{H}(D)$ since $H$ is simply the contraction of the Voronoi partition of $G$. We can now define a graph $G_{D}$ on $N_{H}(D)$ to count the number of cells in $\mathit{resident}(D)$. For every cell $c$ of $\mathit{resident}(D)$, add an edge $e_{c}$ between the balls of $\mathit{hit}(v_{c})\setminus D$. Since the underlying graph $G$ is planar, we can sketch a drawing of $G_{D}$ as follows:
\par
\par
% END S4.SS2.SSS0.P0.SPx3.p3
% BEGIN S4.SS2.SSS0.P0.SPx3.p4
To draw the curve $\zeta_{c}$ corresponding to cell $c$ of $\mathit{resident}(D)$, we first consider the subpath $\pi_{c}$ that we get from Lemma 10, using a plane drawing of $G_{D}$. Let $v$ be the first endpoint of $\pi_{c}$ lying in a ball $D_{1}$ and $v^{\prime}$ be the second endpoint of $\pi_{c}$ lying in a ball $D_{2}$. We first draw the center $x_{D_{1}}$ in its location on any plane drawing of the underlying graph $G$. We then sketch $\zeta_{c}$ on $\pi(x_{D_{1},v})$. This is followed by sketching $\pi_{c}$, and then finally sketching the walk $\pi(v^{\prime},x_{D_{2}})$. Since this sketch is over a path, we will refer to this walk in $G$ as $\pi(\zeta_{c})$.
\par
\par
% END S4.SS2.SSS0.P0.SPx3.p4
% BEGIN S4.F7
\begin{center}\includegraphics[width=0.7\linewidth,height=0.38\textheight,keepaspectratio]{voronoi-page-20.pdf2svg.pdf}\end{center}
\begin{center}\includegraphics[width=0.7\linewidth,height=0.38\textheight,keepaspectratio]{voronoi-page-21.pdf2svg.pdf}\end{center}
  Figure 7: Left: a set of subwalks of cells in $\mathit{resident}(D)$. Right: the walks $\pi(\zeta_{c})$ for each $c\in\mathit{resident}(D)$. 
% END S4.F7
% BEGIN Thmtheorem11
\par\medskip\noindent\textbf{Lemma 11.}\par
For a given ball $D$ in $\mathcal{R}$ and two distinct cells $c$ and $c^{\prime}$ in $\mathit{resident}(D)$, if $\pi(\zeta_{c})$ and $\pi(\zeta_{c^{\prime}})$ share a common vertex $v$, then $e_{c}$ and $e_{c^{\prime}}$ share a common endpoint $D^{\prime}$.
\par
\par
\par
% END Thmtheorem11
% END S4.SS2.SSS0.P0.SPx3
% BEGIN S4.SS2.SSS0.P0.SPx4
\par\medskip\noindent\textbf{Proof.}\par
% BEGIN S4.SS2.SSS0.P0.SPx4.p1
Let $e_{c}=(D_{1},D_{2})$ and $e_{c^{\prime}}=(D_{3},D_{4})$. For $c$ $(c^{\prime})$, the common vertex $v$ must lie on a shortest path from either $x_{D_{1}}$ or $x_{D_{2}}$ ($x_{D_{3}}$ or $x_{D_{4}}$) to $v_{c}$ ($v_{c^{\prime}}$) by construction. Without loss of generality, assume $x_{D_{1}}$ and $x_{D_{3}}$. One of $x_{D_{1}}$ or $x_{D_{3}}$ must be at least as close as the other to $v_{c}$. Without loss of generality, assume $x_{D_{1}}$. By Lemma 8, this implies that $D_{1}$ contains $v_{c^{\prime}}$, and is thus in $\mathit{hit}(v_{c^{\prime}})$. We know that $D$ is contained within $\mathit{hit}(v_{c^{\prime}})$, as we only ever assign $\mathit{base}(v_{c^{\prime}})$ to a member of $\mathit{hit}(v_{c^{\prime}})$. Thus, in order for $c^{\prime}$ to be depth 3, $D_{1}$ must be either the same as $D_{3}$ or $D_{4}$. Thus, $e_{c}$ and $e_{c^{\prime}}$ have a common endpoint. $\hfill\blacktriangleleft$
\par
\par
% END S4.SS2.SSS0.P0.SPx4.p1
% BEGIN S4.SS2.SSS0.P0.SPx4.p2
Lemma 11 implies that $\zeta_{c}$ and $\zeta_{c^{\prime}}$ only cross if $e_{c}$ and $e_{c^{\prime}}$ share an endpoint. Since the only crossings occur between edges incident on the same vertex, $G_{D}$ can be redrawn with these crossings removed. This immediately implies the following:
\par
\par
% END S4.SS2.SSS0.P0.SPx4.p2
% BEGIN Thmtheorem12
\par\medskip\noindent\textbf{Lemma 12.}\par
$G_{D}$ is a simple planar graph.
\par
\par
\par
% END Thmtheorem12
% BEGIN Thmtheorem13
\par\medskip\noindent\textbf{Lemma 13.}\par
There are $O(|\mathcal{R}|)$ cells of depth at most 3.
\par
\par
\par
% END Thmtheorem13
% END S4.SS2.SSS0.P0.SPx4
% BEGIN S4.SS2.SSS0.P0.SPx5
\par\medskip\noindent\textbf{Proof.}\par
% BEGIN S4.SS2.SSS0.P0.SPx5.p1
We first quickly argue that there is $O(n)$ cells of depth at most 2. By the definition of cell, there is at most $n$ cells of depth 1, one for each ball in $\mathcal{R}$. By the existence of a planar dual support (Theorem 2) on $\mathcal{R}$ with respect to all vertices of $G$, we know that for every cell of depth 2 corresponding to vertices $v$ incident on balls $R$ and $R^{\prime}$, the connectivity of $\mathcal{R}_{v}$ requires that $R$ and $R^{\prime}$ share an edge. Since the dual support is planar, there is at most $3|\mathcal{R}|-6$ such pairs, and thus at most $O(n)$ cells of depth 2.
\par
\par
% END S4.SS2.SSS0.P0.SPx5.p1
% BEGIN S4.SS2.SSS0.P0.SPx5.p2
For depth-3 cells, every edge of $G_{D}$ corresponds to a resident depth-3 cell $c$ of $D$, so $|\mathit{resident}(D)|\leq|E(G_{D})|$. By Lemma 12, $|\mathit{resident}(D)|\leq 3|V(G_{D})|$. By charging the depth-3 cells to the total number of edges across all $G_{D}$ for all $D\in\mathcal{R}$ using the planarity of $G_{D}$, $|\mathit{resident}(D)|\leq 3|N_{H}(D)|$. Summing over all neighborhoods in the contracted Voronoi diagram of $\mathcal{R}$ in $G$, we know the sum of degrees is at most $6|\mathcal{R}|$. Thus, $|\mathit{Cell}_{3}|\leq 18|\mathcal{R}|$. $\hfill\blacktriangleleft$
\par
\par
% END S4.SS2.SSS0.P0.SPx5.p2
% END S4.SS2.SSS0.P0.SPx5
% END S4.SS2
% END S4
% BEGIN bib
\par\medskip\noindent\textbf{References}\par
\begin{enumerate}
\item[]  [1]  Brenda S. Baker.  Approximation algorithms for np-complete problems on planar graphs.  Journal of the ACM, 41(1):153–180, January 1994.  doi:10.1145/174644.174650.  
\item[]  [2]  Sayan Bandyapadhyay and Aniket Basu Roy.  Effectiveness of local search for art gallery problems.  In Faith Ellen, Antonina Kolokolova, and Jörg-Rüdiger Sack, editors, Algorithms and Data Structures, pages 49–60, Cham, 2017. Springer International Publishing.  doi:10.1007/978-3-319-62127-2\_5.  
\item[]  [3]  J. Barilan, G. Kortsarz, and D. Peleg.  How to allocate network centers.  Journal of Algorithms, 15(3):385–415, November 1993.  doi:10.1006/jagm.1993.1047.  
\item[]  [4]  Yair Bartal.  Probabilistic approximation of metric spaces and its algorithmic applications.  In Proceedings of 37th Conference on Foundations of Computer Science, pages 184–193, Burlington, VT, USA, 1996. IEEE Comput. Soc. Press.  doi:10.1109/SFCS.1996.548477.  
\item[]  [5]  Aniket Basu Roy.  Covering simple orthogonal polygons with rectangles.  In Alina Ene and Eshan Chattopadhyay, editors, Approximation, Randomization, and Combinatorial Optimization. Algorithms and Techniques (APPROX/RANDOM 2025), volume 353 of Leibniz International Proceedings in Informatics (Lipics), pages 2:1–2:23, Dagstuhl, Germany, 2025. Schloss Dagstuhl – Leibniz-Zentrum für Informatik.  doi:10.4230/LIPIcs.APPROX/RANDOM.2025.2.  
\item[]  [6]  Sujoy Bhore, Balázs Keszegh, Andrey Kupavskii, Hung Le, Alexandre Louvet, Dömötör Pálvölgyi, and Csaba D. Tóth.  Spanners in planar domains via steiner spanners and non-steiner tree covers.  In Proceedings of the 2025 Annual ACM-SIAM Symposium on Discrete Algorithms (SODA), pages 4292–4326. SIAM, 2025.  doi:10.1137/1.9781611978322.145.  
\item[]  [7]  H. Brönnimann and M. T. Goodrich.  Almost optimal set covers in finite vc-dimension.  Discrete \& Computational Geometry, 14(4):463–479, December 1995.  doi:10.1007/BF02570718.  
\item[]  [8]  Kevin Buchin, Marc Van Kreveld, Henk Meijer, Bettina Speckmann, and Kevin Verbeek.  On planar supports for hypergraphs.  In International Symposium on Graph Drawing, pages 345–356. Springer, 2009.  doi:10.1007/978-3-642-11805-0\_33.  
\item[]  [9]  Costas Busch, Ryan LaFortune, and Srikanta Tirthapura.  Sparse covers for planar graphs and graphs that exclude a fixed minor.  Algorithmica, 69(3):658–684, July 2014.  doi:10.1007/s00453-013-9757-4.  
\item[]  [10]  Sergio Cabello and David Gajser.  Simple ptas’s for families of graphs excluding a minor, March 2015.  doi:10.48550/arXiv.1410.5778.  
\item[]  [11]  Timothy M. Chan, Elyot Grant, Jochen Könemann, and Malcolm Sharpe.  Weighted capacitated, priority, and geometric set cover via improved quasi-uniform sampling.  In Proceedings of the Twenty-Third Annual ACM-SIAM Symposium on Discrete Algorithms, pages 1576–1585. Society for Industrial and Applied Mathematics, January 2012.  doi:10.1137/1.9781611973099.125.  
\item[]  [12]  Timothy M. Chan and Sariel Har-Peled.  Approximation algorithms for maximum independent set of pseudo-disks.  Discrete \& Computational Geometry, 48(2):373–392, September 2012.  doi:10.1007/s00454-012-9417-5.  
\item[]  [13]  Hsien-Chih Chang, Jonathan Conroy, Hung Le, Lazar Milenković, Shay Solomon, and Cuong Than.  Covering planar metrics (and beyond): O(1) trees suffice, June 2023.  doi:10.48550/arXiv.2306.06215.  
\item[]  [14]  Hsien-Chih Chang, Jonathan Conroy, Hung Le, Lazar Milenković, Shay Solomon, and Cuong Than.  Shortcut partitions in minor-free graphs: Steiner point removal, distance oracles, tree covers, and more.  In Proceedings of the 2024 Annual ACM-SIAM Symposium on Discrete Algorithms (SODA), pages 5300–5331. SIAM, 2024.  doi:10.1137/1.9781611977912.191.  
\item[]  [15]  M. Chlebík and J. Chlebíková.  Approximation hardness of dominating set problems in bounded degree graphs.  Information and Computation, 206(11):1264–1275, November 2008.  doi:10.1016/j.ic.2008.07.003.  
\item[]  [16]  Kenneth L. Clarkson and Peter W. Shor.  Applications of random sampling in computational geometry, ii.  Discrete \& Computational Geometry, 4:387–421, 1989.  doi:10.1007/BF02187740.  
\item[]  [17]  Kenneth L. Clarkson and Kasturi Varadarajan.  Improved approximation algorithms for geometric set cover.  Discrete \& Computational Geometry, 37(1):43–58, 2007.  doi:10.1007/s00454-006-1273-8.  
\item[]  [18]  Erik D. Demaine, Fedor V. Fomin, Mohammadtaghi Hajiaghayi, and Dimitrios M. Thilikos.  Fixed-parameter algorithms for (k,r)-center in planar graphs and map graphs.  ACM Transactions on Algorithms, 1(1):33–47, July 2005.  doi:10.1145/1077464.1077468.  
\item[]  [19]  Stephane Durocher and Robert Fraser.  Duality for geometric set cover and geometric hitting set problems on pseudodisks.  In CCCG, 2015.  
\item[]  [20]  Zdeněk Dvořák.  Approximation metatheorems for classes with bounded expansion.  LIPIcs, Volume 227, SWAT 2022, 227:22:1–22:17, 2022.  doi:10.4230/LIPIcs.SWAT.2022.22.  
\item[]  [21]  Zdenek Dvorak and Sergey Norin.  Strongly sublinear separators and polynomial expansion, April 2015.  arXiv:1504.04821.  
\item[]  [22]  David Eisenstat, Philip N. Klein, and Claire Mathieu.  Approximating k-center in planar graphs.  In Proceedings of the Twenty-Fifth Annual ACM-SIAM Symposium on Discrete Algorithms, pages 617–627. Society for Industrial and Applied Mathematics, January 2014.  doi:10.1137/1.9781611973402.47.  
\item[]  [23]  David Eppstein.  Diameter and treewidth in minor-closed graph families.  Algorithmica, 27(3):275–291, June 2000.  doi:10.1007/s004530010020.  
\item[]  [24]  Guy Even, Dror Rawitz, and Shimon (Moni) Shahar.  Hitting sets when the vc-dimension is small.  Information Processing Letters, 95(2):358–362, July 2005.  doi:10.1016/j.ipl.2005.03.010.  
\item[]  [25]  Uriel Feige.  A threshold of ln n for approximating set cover.  J. ACM, 45(4):634–652, July 1998.  doi:10.1145/285055.285059.  
\item[]  [26]  Arnold Filtser and Hung Le.  Clan embeddings into trees, and low treewidth graphs.  arXiv:2101.01146 [cs], April 2021.  arXiv:2101.01146.  
\item[]  [27]  Arnold Filtser and Hung Le.  Low treewidth embeddings of planar and minor-free metrics.  arXiv:2203.15627 [cs], March 2022.  doi:10.48550/arXiv.2203.15627.  
\item[]  [28]  Eli Fox-Epstein, Philip N. Klein, and Aaron Schild.  Embedding planar graphs into low-treewidth graphs with applications to efficient approximation schemes for metric problems.  In Proceedings of the 2019 Annual ACM-SIAM Symposium on Discrete Algorithms (SODA), pages 1069–1088. SIAM, January 2019.  doi:10.1137/1.9781611975482.66.  
\item[]  [29]  Greg N. Frederickson.  Fast algorithms for shortest paths in planar graphs, with applications.  SIAM Journal on Computing, 16(6):1004–1022, December 1987.  doi:10.1137/0216064.  
\item[]  [30]  Cyril Gavoille, David Peleg, André Raspaud, and Eric Sopena.  Small k-dominating sets in planar graphs with applications.  In Gerhard Goos, Juris Hartmanis, Jan Van Leeuwen, Andreas Brandstädt, and Van Bang Le, editors, Graph-Theoretic Concepts in Computer Science, volume 2204, pages 201–216. Springer Berlin Heidelberg, Berlin, Heidelberg, 2001.  doi:10.1007/3-540-45477-2\_19.  
\item[]  [31]  Matt Gibson, Gaurav Kanade, Erik Krohn, and Kasturi Varadarajan.  Guarding terrains via local search.  Journal of Computational Geometry, 5(1):168–178, May 2014.  doi:10.20382/jocg.v5i1a9.  
\item[]  [32]  Sathish Govindarajan, Rajiv Raman, Saurabh Ray, and Aniket Basu Roy.  Packing and covering with non-piercing regions.  In Piotr Sankowski and Christos Zaroliagis, editors, 24th Annual European Symposium on Algorithms (ESA 2016), volume 57 of Leibniz International Proceedings in Informatics (Lipics), pages 47:1–47:17, Dagstuhl, Germany, 2016. Schloss Dagstuhl – Leibniz-Zentrum für Informatik.  doi:10.4230/LIPIcs.ESA.2016.47.  
\item[]  [33]  Martin Grohe.  Local tree-width, excluded minors, and approximation algorithms.  Combinatorica, 23(4):613–632, December 2003.  doi:10.1007/s00493-003-0037-9.  
\item[]  [34]  Sariel Har-Peled.  Geometric Approximation Algorithms, volume 173 of Mathematical Surveys and Monographs.  American Mathematical Society, Providence, Rhode Island, June 2011.  doi:10.1090/surv/173.  
\item[]  [35]  Sariel Har-Peled and Kent Quanrud.  Approximation algorithms for polynomial-expansion and low-density graphs.  SIAM Journal on Computing, 46(6):1712–1744, January 2017.  doi:10.1137/16M1079336.  
\item[]  [36]  David Haussler and Emo Welzl.  Epsion-nets and simplex range queries.  Discrete \& Computational Geometry, 2(2):127–151, June 1987.  doi:10.1007/BF02187876.  
\item[]  [37]  D. S. Johnson and H. O. Pollak.  Hypergraph planarity and the complexity of drawing venn diagrams.  Journal of Graph Theory, 11(3):309–325, 1987.  doi:10.1002/jgt.3190110306.  
\item[]  [38]  Gwenaël Joret and Clément Rambaud.  Neighborhood complexity of planar graphs, February 2023.  doi:10.48550/arXiv.2302.12633.  
\item[]  [39]  Michael Kaufmann, Marc van Kreveld, and Bettina Speckmann.  Subdivision drawings of hypergraphs.  In Ioannis G. Tollis and Maurizio Patrignani, editors, Graph Drawing, pages 396–407, Berlin, Heidelberg, 2009. Springer Berlin Heidelberg.  
\item[]  [40]  Klara Kedem, Ron Livne, János Pach, and Micha Sharir.  On the union of jordan regions and collision-free translational motion amidst polygonal obstacles.  Discrete \& Computational Geometry, 1(1):59–71, March 1986.  doi:10.1007/BF02187683.  
\item[]  [41]  Philip Klein, Serge A. Plotkin, and Satish Rao.  Excluded minors, network decomposition, and multicommodity flow.  In Proceedings of the Twenty-Fifth Annual ACM Symposium on Theory of Computing - STOC ’93, pages 682–690, San Diego, California, United States, 1993. ACM Press.  doi:10.1145/167088.167261.  
\item[]  [42]  Hung V. Le.  Structural Results and Approximation Algorithms in Minor-free Graphs.  PhD thesis, Oregon State University, 2018.  
\item[]  [43]  Richard J. Lipton and Robert Endre Tarjan.  A separator theorem for planar graphs.  SIAM Journal on Applied Mathematics, 36(2):177–189, April 1979.  doi:10.1137/0136016.  
\item[]  [44]  Dániel Marx and Michał Pilipczuk.  Optimal parameterized algorithms for planar facility location problems using voronoi diagrams.  In Nikhil Bansal and Irene Finocchi, editors, Algorithms - ESA 2015, pages 865–877, Berlin, Heidelberg, 2015. Springer.  doi:10.1007/978-3-662-48350-3\_72.  
\item[]  [45]  Jir̆í Matous̆ek.  Reporting points in halfspaces.  Computational Geometry: Theory and Applications, 2(3):169–186, November 1992.  doi:10.1016/0925-7721(92)90006-E.  
\item[]  [46]  Jiří Matoušek, Raimund Seidel, and E. Welzl.  How to net a lot with little: Small $\varepsilon$-nets for disks and halfspaces.  In Proceedings of the Sixth Annual Symposium on Computational Geometry, SCG ’90, pages 16–22, New York, NY, USA, May 1990. Association for Computing Machinery.  doi:10.1145/98524.98530.  
\item[]  [47]  Gary L Miller.  Finding small simple cycle separators for 2-connected planar graphs.  Journal of Computer and System Sciences, 32:265–279, 1986.  doi:10.1016/0022-0000(86)90030-9.  
\item[]  [48]  Nabil H. Mustafa and Saurabh Ray.  Improved results on geometric hitting set problems.  Discrete \& Computational Geometry, 44(4):883–895, December 2010.  doi:10.1007/s00454-010-9285-9.  
\item[]  [49]  Nabil H. Mustafa and Kasturi R. Varadarajan.  Epsilon-approximations and epsilon-nets, August 2017.  doi:10.48550/arXiv.1702.03676.  
\item[]  [50]  Jaroslav Nešetřil and Patrice Ossona de Mendez.  Grad and classes with bounded expansion i. decompositions.  European Journal of Combinatorics, 29(3):760–776, April 2008.  doi:10.1016/j.ejc.2006.07.013.  
\item[]  [51]  János Pach and Gerhard Woeginger.  Some new bounds for epsilon-nets.  In Proceedings of the Sixth Annual Symposium on Computational Geometry, SCG ’90, pages 10–15, New York, NY, USA, May 1990. Association for Computing Machinery.  doi:10.1145/98524.98529.  
\item[]  [52]  Ambar Pal, Rajiv Raman, Saurabh Ray, and Karamjeet Singh.  A fast algorithm for computing a planar support for non-piercing rectangles, October 2024.  doi:10.48550/arXiv.2410.02449.  
\item[]  [53]  Evangelia Pyrga and Saurabh Ray.  New existence proofs epsilon-nets.  In Proceedings of the 24th ACM Symposium on Computational Geometry, College Park, MD, USA, June 9-11, 2008, Scg ’08, pages 199–207, New York, NY, USA, 2008. Association for Computing Machinery.  doi:10.1145/1377676.1377708.  
\item[]  [54]  Rajiv Raman and Saurabh Ray.  Constructing planar support for non-piercing regions.  Discrete \& Computational Geometry, 64(3):1098–1122, 2020.  doi:10.1007/s00454-020-00216-w.  
\item[]  [55]  Rajiv Raman and Saurabh Ray.  On the geometric set multicover problem.  Discrete \& Computational Geometry, 68(2):566–591, September 2022.  doi:10.1007/s00454-022-00402-y.  
\item[]  [56]  Rajiv Raman and Karamjeet Singh.  On hypergraph supports.  In European Conference on Combinatorics, Graph Theory and Applications, pages 774–783, 2023.  
\item[]  [57]  Felix Reidl, Fernando Sánchez Villaamil, and Konstantinos Stavropoulos.  Characterising bounded expansion by neighbourhood complexity.  European Journal of Combinatorics, 75:152–168, January 2019.  doi:10.1016/j.ejc.2018.08.001.  
\item[]  [58]  Neil Robertson and P.D Seymour.  Graph minors. xvi. excluding a non-planar graph.  Journal of Combinatorial Theory, Series B, 89(1):43–76, September 2003.  doi:10.1016/S0095-8956(03)00042-X.  
\item[]  [59]  Mikkel Thorup.  Compact oracles for reachability and approximate distances in planar digraphs.  J. ACM, 51(6):993–1024, November 2004.  doi:10.1145/1039488.1039493.  
\item[]  [60]  W. M. van Cleemput.  Hypergraph models for the circuit layout problem.  Applied Mathematical Modelling, 1(3):160–161, 1976.  
\item[]  [61]  Vladimir N Vapnik and A Ya Chervonenkis.  On the uniform convergence of relative frequencies of events to their probabilities, 1971.  
\item[]  [62]  Kasturi Varadarajan.  Epsilon nets and union complexity.  In Proceedings of the Twenty-Fifth Annual Symposium on Computational Geometry, pages 11–16, Aarhus Denmark, June 2009. ACM.  doi:10.1145/1542362.1542366.  
\item[]  [63]  Kasturi Varadarajan.  Weighted geometric set cover via quasi-uniform sampling.  In Proceedings of the Forty-Second ACM Symposium on Theory of Computing, pages 641–648, Cambridge Massachusetts USA, June 2010. ACM.  doi:10.1145/1806689.1806777.  
\item[]  [64]  A. A. Voloshina and V. Z. Feinberg.  Planarity of hypergraphs.  In Doklady akademii nauk belarusi, volume 28, pages 309–311, Minsk, Belarus, 1984. Academii Nauk Belarusi.  
\item[]  [65]  Christian Wulff-Nilsen.  Separator theorems for minor-free and shallow minor-free graphs with applications.  In 2011 IEEE 52nd Annual Symposium on Foundations of Computer Science, pages 37–46, 2011.  doi:10.1109/FOCS.2011.15.  
\end{enumerate}
% END bib
